\paragraph{Spin term, Lichnerowicz contact and the smooth first-order cone in the comparison of \(G_T\) with \(G_A\) (conditional; not closure).}
Keep the setting of the paragraphs on fermionic content and on fermion replica gluing: Poisson limit, Euclidean section, \(\alpha=1\), the \(\LUV^2\) term, \(N_s\) scalars, the connection with its ghosts and \(n_D\) Dirac fields, \(N_{\mathrm{tot}}=N_s+2-4n_D\), and the coframe \(e^A{}_\mu=\eta^{AB}\omega_{\mu B}\) with torsion \(T\) of the paragraph on the torsion invariant. The \emph{fixed-state replica} is the pulled-back replica of the paragraph on the replica cone coefficient (\(\rho\) regular, cone on \(G\)) with the antiperiodic Dirac gluing of the paragraph on fermion replica gluing. At random two-jets of a closed \(\omega\), the operator \(\mathcal O_f\) of the fermionic-content paragraph satisfies \(\tfrac14\operatorname{tr}E_f=\tfrac14R_G-\tfrac1{16}T_{\lambda\mu\nu}T^{\lambda\mu\nu}+\tfrac18T_{\lambda\mu\nu}T^{\mu\lambda\nu}\), with no further term in the representation by \(R_G\), \(T^2\) and \(TT'\); the scalar identity differs only in the sign of the first torsion term, so that \(Q_f=Q_s=\tfrac14\) in the torsion reading, whereas in the variables \((G,\rho,\partial\Theta)\) one has \(Q_f=0\). The coefficient of \(R_G\) in \(\operatorname{tr}(\tfrac16R_G-E_f)\) is then \(-\tfrac13\) over four components, and with weight \(-1\) it gives \(+\tfrac13\) per Dirac field, the antiperiodic cone value, instead of \(-\tfrac23\). In the normalisation of the paragraph on fermion replica gluing, the \(n\)-derivative at \(n=1\) of the heat-trace constant of a two-component block is \(\tfrac16\) for the antiperiodic gluing with the less singular extension and \(-\tfrac13\) for the unrotated gluing with the componentwise-regular domain; the difference \(\tfrac12=2Q_f\) equals the contact term of \(Q_fR_G\) continued with its conical delta. This is an equality of \(n\)-derivatives at \(n=1\), not an identity in \(n\). The excess \(+1\) per Dirac field of the antiperiodic over the unrotated area coefficient is therefore the Lichnerowicz contact. It is already contained in \(c_A\), it enters the area side only and raises \(c_A\), and adding to it the contact of \(Q_fR_G\), which at the weight \(-1\) shifts \(c_A\) by \(+1\) per Dirac field, would count it twice. For a Dirac field with the Levi-Civita spin connection the black-hole and entanglement entropies coincide~\cite{Kabat1995}; for the torsion reading the corresponding equality (both \(+\tfrac13\) per Dirac field) is the computation above.

\emph{Smooth first-order cone.} For \(\omega=\bar\omega+\epsilon\,\mathrm dA\), with any constant antisymmetric invertible \(\bar\omega\) and any \(A\), one has to first order in \(\epsilon\) \(R_G^{(1)}=-2\Box_{\bar G}\,\delta\ln\rho\) and \(\bar\rho\,R[g_\rho]^{(1)}=-\tfrac12\bar\rho^{\,2}R_G^{(1)}\). These were checked symbolically for a generic \(\bar\omega\) and follow from the Weitzenb\"ock identities of the paragraph on the torsion invariant. A conical deficit \(2\pi(1-n)\) at first order in \(n-1\) of a smooth Poisson structure about constant \(\bar\Theta\) therefore requires \(\ln\rho\) to carry the flux \(\oint\nabla\ln\rho\) through \(\Sigma\), that is a singular \(\rho\); the fixed-state replica is not a limit of smooth Poisson structures at this order. This does not by itself exclude the fixed-state replica as the definition of \(\operatorname{Tr}\rho^n\). For a regularised smooth family the first-order local density of every sector has \(Q=\tfrac14\), the factor \(-\tfrac12\) of the paragraph on the replica cone coefficient then carries no branch dependence, and the area coefficient is \(c_A^{\mathrm{sm}}=-N_{\mathrm{tot}}/12\), the factor \(-\tfrac12\) times the coefficient of \(\rho R[g_\rho]\) in every sector including the Dirac fields, for which it coincides with the antiperiodic value \(+\tfrac13\) without any gluing prescription. The branch issue of the unregularised singular-\(\rho\) cone is not removed. Thus \(c_A^{\mathrm{sm}}\) equals the fixed-state coefficient \(c_A\) plus the bosonic contacts \(-\tfrac14(N_s+2)\), and the two readings differ by the total derivative \(\Box\delta\ln\rho\), that is by the boundary flux of \(\nabla\ln\rho\) through \(\Sigma\).

\emph{Outcome for \(n_D=2\), \(N_s=6\).} Here \(c_T=0\) and \(N_{\mathrm{tot}}=0\). The fixed-state replica gives \(c_A=2\). If the connection sector carries a contact of the Kabat value \(-1\) in these units (the value that takes the vector with its ghosts from \(\tfrac13\) to \(-\tfrac23\)), which for a Maxwell field is the edge-mode term of Donnelly and Wall~\cite{DonnellyWall2016}, then \(c_A=1\). The smooth reading gives \(c_A^{\mathrm{sm}}=0=c_T\), an equality of vanishing coefficients, which holds for every sector-uniform \(Q\) because \(N_{\mathrm{tot}}=0\) and is not a derivation. For fields minimally coupled in the variables \((G,\rho,\partial\Theta)\), the contacts of these types have weight \(+1\) per Dirac field (already contained in \(c_A\)), \(-1\) for the vector and \(0\) for the scalars and ghosts, so within this class they give at most the value \(1\) above; the torsion-reading bosonic contacts \(-\tfrac14(N_s+2)\) are the smooth reading and do not belong to this class. The values \(2\) and \(0\) correspond to a replica in which \(\ln\rho\) carries no flux through \(\Sigma\) and one in which it carries the full deficit; intermediate fluxes are not computed, and which replica the matrix state selects depends on the state and the tip regularisation, which are not derived here. These statements are for the Poisson-limit operators, the \(\LUV^2\) term, one loop, \(\alpha=1\) and the fixed-state and first-order smooth replicas stated; they do not derive the replica, the state, the tip domain or the gluing, do not treat edge modes of the symplectic gauge sector of \(a^\mu\), of the ghosts or of fermions (not examined), the \(\LUV^4\) term for \(N_{\mathrm{tot}}\neq0\), or the next order of either side, and give no value or sign of \(G_{\mathrm{eff}}\). Problems~B, C, D, F and~G remain open.
